\section{From Quantum Memories to EIT Storage}
\label{sec:from_quantum_memories_to_eit}

Photons are natural carriers of quantum information over long distances. Their weak interaction with the environment allows optical quantum states to be transmitted through fibres or free-space links with low noise. The same weak interaction, however, makes deterministic storage challenging, since a freely propagating photon must be absorbed, delayed, or coherently converted into another degree of freedom before it can be held. Quantum networks therefore require a reversible interface between travelling optical states and stationary matter excitations~\cite{AWAVQMitCoSR_Rodriguez_,Oqmboeit_MaEtAl_2017}.

An optical quantum memory performs this interface by writing an input optical state into a material system, preserving it for a chosen storage time, and retrieving it into a well-defined optical mode. Following the standard description of optical quantum memories~\cite{AWAVQMitCoSR_Rodriguez_,Oqmboeit_MaEtAl_2017}, the idealized process may be represented as
\begin{equation}
    \ket{\psi}_{\mathrm{light}}
    \ket{G}_{\mathrm{atoms}}
    \longrightarrow
    \ket{0}_{\mathrm{light}}
    \ket{\psi}_{\mathrm{atoms}}
    \longrightarrow
    \ket{\psi}_{\mathrm{light}},
    \label{eq:quantum_memory_mapping}
\end{equation}
where $\ket{\psi}_{\mathrm{light}}$ denotes the incoming optical state, $\ket{G}_{\mathrm{atoms}}$ denotes the initial atomic state, and $\ket{\psi}_{\mathrm{atoms}}$ denotes the stored matter excitation. The performance of the mapping is determined by efficiency, fidelity, bandwidth, noise, storage time, and multimode capacity, with different physical platforms optimizing different combinations of these quantities~\cite{AWAVQMitCoSR_Rodriguez_,Oqmboeit_MaEtAl_2017}.

A single isolated atom contains internal states that could, in principle, store quantum information, but its coupling to a freely propagating optical mode is weak. After optical excitation, spontaneous emission generally distributes the excitation over many modes rather than returning it coherently to the prescribed one. Without a cavity or another enhancement mechanism, efficient absorption and re-emission into a selected spatial mode are therefore strongly suppressed~\cite{AWAVQMitCoSR_Rodriguez_}. Atomic ensembles overcome this limitation through collective coupling. The optical excitation is delocalized over many atoms, and the coherent addition of the emission amplitudes enhances the coupling between the optical field and the stored matter coherence~\cite{AWAVQMitCoSR_Rodriguez_,Oqmboeit_MaEtAl_2017}.

The stored excitation is described by a collective spin wave. If all atoms are initially prepared in the ground state $\ket{g}$, the absorption of a single photon can transfer one atom to a long-lived state $\ket{s}$, where no measurement reveals which atom has changed. Following the standard collective-excitation description of ensemble quantum memories~\cite{AWAVQMitCoSR_Rodriguez_,Oqmboeit_MaEtAl_2017}, the single-excitation spin wave is written as
\begin{equation}
    \ket{1_S}
    =
    \frac{1}{\sqrt{N}}
    \sum_{j=1}^{N}
    e^{i\Delta\mathbf{k}\cdot\mathbf{r}_j}
    \ket{g_1,\ldots,s_j,\ldots,g_N},
    \label{eq:intro_spin_wave}
\end{equation}
where $N$ is the number of atoms, $\mathbf{r}_j$ is the position of atom $j$, and $\Delta\mathbf{k}$ is the wave-vector mismatch transferred from the optical fields to the atomic coherence. The phase factor in Eq.~\eqref{eq:intro_spin_wave} determines the spatial grating of the stored coherence and therefore fixes the phase-matching condition for retrieval. Preservation of the collective phase is consequently required for conversion back into the desired optical mode~\cite{AWAVQMitCoSR_Rodriguez_,Oqmboeit_MaEtAl_2017}.

A two-level atom is not suffcient to provide a suitable long-lived storage state for this mapping. If the optical field couples only $\ket{g}$ to an excited state $\ket{e}$, the excitation remains stored in an optically excited state that is rapidly lost by spontaneous emission. A third state $\ket{s}$ is required, with a lifetime much longer than that of $\ket{e}$ and weak coupling to the electromagnetic vacuum. The optical field can then be mapped onto a ground-state coherence between $\ket{g}$ and $\ket{s}$ rather than onto population in the short-lived excited state~\cite{AWAVQMitCoSR_Rodriguez_,Oqmboeit_MaEtAl_2017}.

Ground-state optical storage is therefore naturally described by three-level light-matter systems. Depending on the ordering of the levels one obtains $\Lambda-$, ladder, or $V-$types system. For atomic quantum memories, the relevant configuration is usually the $\Lambda$ system, in which two long-lived lower states, $\ket{g}$ and $\ket{s}$, are optically coupled to a common excited state $\ket{e}$. Where one optical field, the signal field, couples the transition $\ket{g}\leftrightarrow\ket{e}$, while a stronger control field couples $\ket{s}\leftrightarrow\ket{e}$. The control field determines how the optical excitation is transferred into the ground-state coherence and how it is later converted back into light~\cite{Oqmboeit_MaEtAl_2017,Apgteitiav_FinkelsteinEtAl_2023}.

Electromagnetically induced transparency (EIT) provides a controlled storage protocol for this $\Lambda$ system. In EIT, the control field modifies the optical susceptibility of an otherwise absorbing ensemble. A resonant signal field can propagate through a narrow transparency window created by quantum interference, while the associated steep dispersion reduces the group velocity and spatially compresses the pulse inside the medium. In the standard dark-state polariton treatment~\cite{C:Tampsiae_Lukin_2003}, adiabatic reduction of the control field converts the optical component of the excitation into a ground-state atomic coherence. Reapplying the control field reverses the mapping and retrieves the stored optical field~\cite{Oqmboeit_MaEtAl_2017,Apgteitiav_FinkelsteinEtAl_2023,C:Tampsiae_Lukin_2003}.

Warm alkali vapours provide a practical platform for EIT storage because laser cooling and cryogenic operation are not required, while the optical depth can be controlled through the vapour temperature. Cesium (Cs) and Rubidium (Rb) are commonly used in this context because their D-line transitions are accessible with near-infrared lasers and their hyperfine ground states can support long-lived coherences. The compact quantum memory demonstrated by Jutisz \emph{et al.} is particularly relevant here, since it implements EIT storage in a Cs vapour cell using the D$_1$ line and provides the experimental reference system for part of the modelling developed later~\cite{Smqms_JutiszEtAl_2025}. In such warm-vapour memories, the retrieval efficiency is not determined by the internal level structure alone. Atomic motion, wall collisions, magnetic dephasing, and finite spatial mode overlap all modify the stored spin wave and must be included when the retrieved optical field is modelled~\cite{Smqms_JutiszEtAl_2025,Apgteitiav_FinkelsteinEtAl_2023}.

The remaining discussion follows the physical sequence required to describe EIT storage in a warm-vapour ensemble. The alkali-metal level structure is first specified, since the hyperfine ground states provide the two long-lived states of the effective $\Lambda$ system and the excited D-line manifold supplies the optically coupled state. The EIT mechanism is then formulated in terms of the signal and control fields, which determine how the travelling optical excitation is converted into a collective ground-state coherence. After storage, the relevant dynamical object is the spin wave. Its amplitude envelope, phase factor, and subsequent evolution under atomic motion determine the spatial mode and efficiency of the retrieved field.

